% \begin{abstract}
% Industrial anomaly detection is critical for product quality assurance and process reliability, yet practical inspection requires more than determining whether an image is normal or anomalous. A useful system should also localize defective regions, recognize different defect types, and provide interpretable defect analysis and diagnosis. Existing CLIP-based methods have shown strong performance in binary anomaly detection and segmentation, yet most of them remain limited to binary discrimination and offer only weak support for fine-grained defect understanding. Meanwhile, recent autoregressive vision-language models (VLMs) provide powerful multimodal reasoning and defect description capabilities, but existing VLM-based approaches are limited to text outputs and cannot directly produce pixel-level segmentation masks. As a result, a generalist model that unifies detection, segmentation, multi-type anomaly recognition, and defect analysis remains missing. 
% We present GenAU, a \textbf{G}eneralist vision-language framework for industrial \textbf{A}nomaly \textbf{U}nderstanding that unifies anomaly detection, anomaly segmentation, multi-type anomaly detection, and defect analysis within a single model. Specifically, we extend an instruction-following VLM with two segmentation tokens, [SEG DEFECT] and [SEG NORMAL], and align their embeddings with multi-scale visual features from a frozen vision encoder via adapters to generate fine-grained anomaly maps. At the same time, the language head predicts defect-aware textual outputs, enabling the model to recognize and describe specific defect types in natural language. This design allows GenAU to combine the reasoning strengths of large VLMs with pixel-level localization, moving beyond text-only defect recognition toward joint semantic understanding and spatial grounding.
% \end{abstract}

\begin{abstract}
Industrial inspection requires more than binary anomaly detection: a practical system should determine whether an anomaly exists, localize the defective region, identify the defect type, and provide interpretable visual evidence for its decision. Existing CLIP-based methods show strong anomaly detection and localization but offer limited support for language-level defect understanding, while instruction-tuned vision-language models can describe defects but do not natively produce pixel-level masks. We introduce GenAU, a \textbf{Gen}eralist vision-language framework for industrial \textbf{A}nomaly \textbf{U}nderstanding that unifies image-level anomaly detection, pixel-level anomaly segmentation, multi-type anomaly detection, and defect analysis within a single instruction-following architecture. GenAU augments a vision-language model with two segmentation tokens, \texttt{[SEG\_defect]} and \texttt{[SEG\_normal]}, whose hidden states serve as language-grounded queries for multi-scale visual features, enabling pixel-level anomaly localization. The image-level anomaly score fuses this localization map with the decoder's textual normal/defect decision. The language decoder generates structured defect-aware responses, allowing GenAU to recognize multiple anomaly types and provide textual analysis in natural language. We train GenAU with a joint objective combining language modeling and pixel-level segmentation. Experiments under cross-dataset industrial anomaly benchmarks show that GenAU achieves the strongest image-level anomaly detection among CLIP-based zero-shot methods on VisA and Real-IAD and the strongest segmentation on MPDD, while its segmentation on VisA and its image-level AUROC on MPDD's logical, large-area anomalies stay below the best CLIP-based baselines. The same checkpoint covers all four tasks, adding zero-shot multi-type detection and language-grounded defect analysis.
\end{abstract}
